\section{Discussion}

\subsection{Key Findings and Implications}

\textbf{Contamination-inflation correlation is measurable.} Our primary finding---Spearman $r = 0.326$---demonstrates that contamination effects are systematic enough to model. This establishes a foundation for contamination-aware evaluation.

\textbf{Magnitude is moderate, not large.} The observed correlation is weaker than our primary target ($r > 0.5$), suggesting that while contamination inflates scores, capability remains the primary driver of benchmark performance.

\subsection{Limitations}

\textbf{Single model family.} We analyze only Pythia models trained on The Pile. Generalization to other architectures requires additional validation.

\textbf{Corpus coverage.} Our 50,000 document subset represents 0.006\% of The Pile. Full corpus analysis would provide definitive overlap statistics.

\textbf{Simulated validation.} Current results derive from proof-of-concept validation using simulated contamination data where contamination-inflation correlation was structurally embedded in the data generation process. Full experimental validation requires running actual Pythia checkpoint evaluations with a complete Pile n-gram index ($\sim$144 GPU-hours for evaluation, $\sim$48 CPU-hours for index construction). The correlation statistics presented demonstrate methodological feasibility; real Pythia checkpoint experiments are required to validate whether the observed relationship holds in practice.

\textbf{Untested causal chain.} We establish correlation but do not verify the full causal mechanism (exposure $\rightarrow$ memorization $\rightarrow$ correct answers).

\textbf{Statistical significance caveat.} Only 2 of 4 benchmarks (MMLU, ARC-Challenge) show statistically significant correlation at $p < 0.05$. HellaSwag ($p = 0.08$) and WinoGrande ($p = 0.15$) do not reach significance.

\subsection{Broader Impact}

Knowledge of contamination-inflation relationships enables contamination-aware model comparison. The same detection methods used in this work enable defense against adversarial training that maximizes benchmark scores through contamination rather than capability.
